\chapter{Universal System Exclusive}

\section{What is recognised}
\estab\ The instrument recognises exactly \textbf{two} universal messages:

\begin{center}
\begin{tabular}{ll}
\toprule
bytes & meaning \\
\midrule
\bytes{F0 7E 7F 09 01 F7} & GM System On \\
\bytes{F0 7E 7F 09 02 F7} & GM System Off \\
\bottomrule
\end{tabular}
\end{center}

\estab\ Both are also \emph{transmitted} by the instrument, from literals at
\bytes{0xF4FEE6} and \bytes{0xF4FEEC}, and both are written into Standard MIDI Files as a
delta-time-zero five-byte SysEx event (\bytes{00 F0 05 7E 7F 09 01 F7}).

\estab\ GM System Off is a no-op when GM mode is already off.

\section{What is not recognised}
\estab\ The following are rejected by the grammar. A device expecting any of them will get
no response:

\begin{itemize}
\item \bytes{F0 7F 7F 04 01 \textit{ll} \textit{mm} F7} --- Master Volume. The whole
  \bytes{0x7F} Universal \emph{Real Time} identifier is unmatched at byte~1.
\item \bytes{F0 7E \textit{dd} 06 01 F7} --- Identity Request. There is no identity reply
  in this firmware.
\item \bytes{F0 7E 7F 09 03 F7} --- GM2 System On.
\end{itemize}

\section{Two caveats}
\likely\ Trailing bytes before the \bytes{F7} are probably tolerated, because the grammar
stops matching once a command is resolved.

\likely\ The identifier byte is not re-checked after the front-end accepted it, so
\bytes{F0 50 7F 09 01 F7} would probably act as GM System On despite carrying the
Matsushita identifier. Neither caveat was executed on hardware or in emulation.
